\noindent\begin{minipage}{\textwidth}
\begin{center}
\thispagestyle{empty}
\vspace{0.5cm}
{ \Large{Prospects of Color Transparency Measurements with the
SPD Detector at the NICA Collider}}\\
\vspace{1cm}

{\large Iu.A. Skorodumina$^{1, a}$ and A.B. Larionov$^{2}$} \\[16pt]

\parbox{.86\textwidth}{\centering\footnotesize\it
$^1$DLNP, Joint Institute for Nuclear Research, 141980 Dubna, Russia\\[3pt]
\setstretch{0.3} 
$^2$BLTP, Joint Institute for Nuclear Research, 141980 Dubna, Russia\\[18pt]
%\setstretch{0.3} 

E-mail: $^a$skorodumina@jinr.ru}\\


\vspace{2cm}
{\bf Abstract}\\[9pt]

\end{center}
{\small We suggest to measure color transparency effects for the process $pd\rightarrow ppn$ (proton knock-out from a deuteron) in $d-d$ collisions with the SPD detector at the NICA collider. This report thoroughly investigates the overall feasibility of such measurements and describes the general measurement strategy. Observables sensitive to color transparency effects are discussed. The Monte Carlo simulation of the process $pd\rightarrow ppn$ was performed based on the theoretical cross section approximation obtained in the framework of the generalized eikonal approach supplemented by the quantum diffusion model of the color transparency. Estimations of the event rates for the proposed experiment are presented. Favorable kinematics for the proposed measurements is discussed. The specificities of the proton and neutron detection during the first stage of the SPD experiment are described. A preliminary event selection procedure for the proposed measurements was developed. The main challenges and potential roadblocks of the proposed measurements are addressed, as well as possible solutions to them.}


\end{minipage}

